\begin{frame}{교차 엔트로피를 손실함수로: 직관}
  \begin{itemize}
    \item 로지스틱회귀에서 정답 $y^{(i)}$ = ``정답 분포''
      (클래스 1 확률 100\% 또는 0\%), $h_{\boldsymbol{w}}(\boldsymbol{x}^{(i)})$ = 모델이 예측한
      분포.
    \item \textbf{교차 엔트로피를 최소화하는 것 = 모델의 예측
      분포를 정답 분포에 최대한 가깝게 만드는 것} --- 손실함수로
      쓰는 이유:
      \[
      J(\boldsymbol{w}) = -\frac{1}{m}\sum_{i=1}^m \left[y^{(i)}\log h_{\boldsymbol{w}}(\boldsymbol{x}^{(i)}) +
      (1-y^{(i)}) \log(1-h_{\boldsymbol{w}}(\boldsymbol{x}^{(i)}))\right]
      \]
    \item ML에서는 보통 $\log_2$ 대신 \textbf{자연로그} $\ln$ ---
      상수배($1/\ln 2$)만 차이.
  \end{itemize}
  \vskip0.4em
  {\footnotesize 직관: 정답 $y = 1$ 인데 $h_{\boldsymbol{w}}(\boldsymbol{x}) \to 0$ 으로 확신하면
  $-\log h_{\boldsymbol{w}}(\boldsymbol{x}) \to \infty$ --- \textbf{확신을 갖고 틀리면 그만큼
  크게 벌한다}.}
\end{frame}
